\section{Introduction}\label{sec:intro}

Regenerative chatter limits the material removal rate of most milling operations. Its linear mechanism, the delayed feedback of the surface left by the previous tooth, is captured by stability lobe diagrams computed in the frequency domain with the zero-order approximation \citep{altintas1995analytical} or in the time domain with semi-discretisation and related methods \citep{insperger2002semi,insperger2004updated,insperger2011book}; see the reviews of \citet{quintana2011chatter}, \citet{munoa2016chatter} and \citet{altintas2020chatter}. Among the suppression techniques, active damping adds a controlled force to the machine or the workpiece. Proof-mass (inertial) actuators driven by direct velocity feedback (DVF) are attractive because the loop is simple and collocated: with an ideal force actuator it is unconditionally stable, and with an inertial actuator it is stable up to a gain set by the actuator's own resonance \citep{elliott2001feedback}; active damping of milling chatter with such actuators and with related structural actuators has been demonstrated on laboratory set-ups and machine tools \citep{dohner2004mitigation,munoa2013chatter,vandijk2012robust,monnin2014optimal}.

The weak point of an inertial actuator is the force it can deliver. A stroke-limited inertial actuator in a velocity feedback loop damps small vibrations and can lose stability at large ones, which gives rise to limit cycles \citep{baumann2007destabilization}. In milling the actuator has to provide two forces at once: the force that damps the chatter mode, and the force that the feedback reacts with against the forced, tooth-periodic vibration of a stable cut. \citet{ozsoy2025frequency} modelled actuator saturation in the frequency domain, and \citet{ozsoy2026robust} proposed a saturation-aware gain map: at each spindle speed, the DVF gain that maximises the linear stability limit subject to the peak force demand of a simulated cut staying within the actuator capacity. Their study is the starting point of the present one.

Saturation turns the controlled cut into a nonlinear system, and nonlinear chatter has a long history. Loss of contact between tool and workpiece bounds the chatter amplitude \citep{tlusty1981basic}, the Hopf bifurcation of the regenerative model is subcritical \citep{kalmarnagy2001subcritical,stepan2005nonlinear}, and linearly stable cuts coexist with chatter in bistable, or unsafe, zones below the linear limit \citep{dombovari2015bistable}. Bistability of turning under active velocity feedback damping has been analysed by \citet{sahu2021bistability}. For chatter controllers designed with input constraints, regional guarantees around the chatter-free motion have been given as ellipsoidal invariant sets of a predictive controller \citep{zhang2015ellipsoid} and as positively invariant sets of an adaptive controller with hard input clipping \citep{wu2016adaptive}. The tools for such guarantees come from the analysis of linear systems with saturating inputs: sector conditions that bound the dead-zone nonlinearity locally and give ellipsoidal estimates of the region of attraction through linear matrix inequalities (LMIs) \citep{hu2001book,gomesdasilva2005antiwindup,tarbouriech2011book}, extended to periodic systems with periodic or clock-dependent Lyapunov functions \citep{allerhand2011robust,briat2013clock,li2015periodic,zhou2011periodic,yang2015periodic}. Delay equations are commonly turned into ordinary differential equations of finite order by spectral collocation of the delayed history \citep{breda2005pseudospectral,breda2015book}, an approach also used for the periodic milling equation \citep{butcher2009chebyshev}.

A second line of work reduces the control effort itself. A delayed feedback of the form $\dot x(t)-\dot x(t-T)$ vanishes on every motion with period $T$ \citep{pyragas1992continuous}; in machining, controllers built on the difference between the present and the delayed motion have been used to shape the stability limit \citep{vandewouw2015ifac,vandijk2015ijrnc}, to counter the regenerative effect in turning \citep{mancisidor2019mechatronics}, and to obtain stability with less control effort \citep{li2022ijamt,du2024ijms}. Such a term would spend none of the actuator's authority on the forced vibration of a stable cut, which is exactly the share that a saturation-aware DVF design has to budget.

What remains open, for a fixed DVF law whose command is clipped at the actuator capacity in a regenerative milling model, is (i) how the authority condition of a saturation-aware gain map should be posed so that it is well defined, (ii) whether and by how much the clip makes the linearly stable cut bistable across the speed range, and how this depends on the way the tool enters the cut, (iii) how large a basin around the forced orbit can be guaranteed when part of the authority is spent on the forced vibration, and (iv) whether a delayed term that spends no authority on the forced orbit enlarges the basin, and at what cost in robustness. This paper addresses these four questions on the case study of \citet{ozsoy2026robust}. Its contributions are:
\begin{enumerate}
\item A well-posed authority condition (Section~\ref{sec:feasibility}). The published test takes the peak demand of a simulation that starts from rest, which is set by the entry transient and makes the feasible set of gains fragmented. The condition proposed here bounds the steady forced demand in closed form at every depth up to a margin below the limit and checks the result in the time domain with the clipped actuator.
\item A map of saturation-induced bistability for the clipped DVF loop over 251 spindle speeds (Section~\ref{sec:bistability}): the dependence of the attainable depth on the entry, the lower edge of the bistable band, and the size of the basin measured by velocity kicks, with and without the clip, in a simulator that remembers the full machined surface.
\item Certificates of the basin (Section~\ref{sec:certificates}): ellipsoidal sets around the forced orbit proved by sector-condition LMIs in which the dead zone moves with the forced command, so that the authority spent on the forced vibration enters the certificate, first on a collocation model of the averaged milling model and then on one of the time-periodic model, each solution checked independently of the solver and validated against the simulator.
\item An analysis of the authority-preserving delayed term under the same hard clip and the same authority rule as DVF (Section~\ref{sec:delayed}): linear limits, simulated entries, certified basins, and the robustness limits set by loop delay, unmodelled modes, delay mismatch and runout.
\end{enumerate}
The tools are not new: the bistable zone, sector conditions, collocation and the delayed difference law are all established. The contribution is their combination for this problem, the quantification over the speed range, and the cross-validation of the certificates against an independent simulator. The study is computational; Section~\ref{sec:discussion} discusses what an experimental validation would need. Section~\ref{sec:model} introduces the model and the case study, and Section~\ref{sec:conclusions} concludes.
